\documentclass[10pt,a4paper]{amsart}
\usepackage[margin=19mm]{geometry}
\usepackage{amsmath,amssymb,mathtools}
\usepackage[scr=rsfs,cal=euler]{mathalfa}
\usepackage{xfrac}
\usepackage[unicode,hidelinks,bookmarks=false]{hyperref}
\newcommand{\ie}{i.e.\ }
\newcommand{\cf}{cf.\ }
\newcommand{\resp}{respectively\ }
\newcommand{\Subsection}{\subsection}
\providecommand{\nobreakdash}{\nobreak\hbox{-}}
\setlength{\emergencystretch}{2em}
\multlinegap0pt
\usepackage{amssymb}
\usepackage{dsfont}
\usepackage[shortlabels]{enumitem}
\usepackage{float}
\newtheorem{theorem}{Theorem}
\title{Fractional differ-integral involving bicomplex Prabhakar function in the kernel and applications}
\author{Urvashi Purohit Sharma, Ritu Agarwal}
\date{Source: arXiv:2603.07034v1 (CC BY 4.0)}
\begin{document}
\maketitle

\noindent\emph{Source and licence.} Fractional differ-integral involving bicomplex Prabhakar function in the kernel and applications. Version arXiv:2603.07034v1 (CC BY 4.0), distributed by arXiv under the Creative Commons Attribution 4.0 International licence, http://creativecommons.org/licenses/by/4.0/, as recorded in the arXiv licence metadata for this exact version and shown on the abstract page https://arxiv.org/abs/2603.07034v1. Full licence notice: \texttt{licenses/arxiv-2603.07034-CC-BY-4.0.txt}.\par\smallskip

\noindent\textit{Editorial scope.} Counted unit: the article's theorem th:cauchy-nonhomo (source0.tex lines 955--974). The article's complete original proof of it is delivered with it (source0.tex lines 975--1080, one \texttt{proof} environment of 106 lines). No mathematics is written, repaired or re-derived: the statement and the proof are the article's own lines, in the article's own notation, and no retained line is substituted or re-flowed.  No further source-local context is needed: every cross-reference of every retained line resolves inside the two delivered spans.  Nothing was omitted from any retained span, so the package carries no omission record.  No retained line contains a citation, so the wrapper carries no bibliography and no bibliography entry.  No retained line contains a figure environment, an image-inclusion command or drawing code, so no image is delivered and the package declares no asset dependency.  Editorial formatting only, in addition: an AMS \texttt{amsart} preamble that restates the article's own declarations and macros used by the retained lines, namely the environments \texttt{theorem} on separate counters, together with the standard packages those lines need.  The retained environments print in this document as Theorem 1 in order of appearance, whereas the article prints the same items with the numbering of its own full document.  The complete native source of the article as distributed in the arXiv e-print for this exact version is bundled unchanged as \texttt{sources/195912/source0.tex}.
\begin{theorem}\label{th:cauchy-nonhomo}
    Let $f,g\in  L_{\mathbb{T}}[a, b] $ and of exponential order, be bicomplex valued functions and
\begin{equation}\label{cauchy1}
{}^{C}\mathbb{D}^{l}_{m,n,r,0^+} f(t)+kf(t) = g(t), 
\qquad f(0^+) = \tau_0,
\end{equation}
where $k$ and $\tau_0$ are constants and the parameters satisfy 
$m,n,r \in \mathbb{T},~|\operatorname{Im_j}(m)|<\operatorname{Re}(m),~|\operatorname{Im_j}(n)|<\operatorname{Re}(n)$ and $l\in \mathbb{D},~l=l_1e_1+l_2e_2, ~0<l_1,l_2<1.$


The integral is taken from the lower limit $0^+$. Then
\begin{equation}
f(t)= (h *g)(t)+\tau_{0}\mathcal{L}^{-1}\Big\{\frac{\xi^{n-1}(1 - r \xi^{-m})^{l}}{\xi^{n}(1 - r \xi^{-m})^{l} + k}\Big\}(t),
\end{equation}
where 
%\[E_{\alpha}(-k t^{\alpha})+ \int_{0}^{t} (t - \tau)^{\alpha - 1}E_{\alpha, \alpha}\!\left(-k (t - \tau)^{\alpha}\right)g(\tau)\, d\tau\]
\[
h(t) = \mathcal{L}^{-1}\left\{\frac{1}{\xi^{n}\big(1 - r\xi^{-m}\big)^{l} + k}\right\}(t),\]

\end{theorem}	

    \begin{proof}
        	Taking the Laplace transform of \eqref{cauchy1} and using the formula  
	\[
	\mathcal{L}({}^{C}\mathbb{D}^{l}_{m,n,r,0^+} f(t) ;\xi)
	= \xi^{n}(1 - r \xi^{-m})^{l}\tilde{f}(\xi) - \xi^{n-1}(1 - r \xi^{-m})^{l}f(0^+),
	\]
	we obtain
	\[
	\xi^{n}(1 - r \xi^{-m})^{l}\tilde{f}(\xi) - \xi^{n-1}(1 - r \xi^{-m})^{l}\tau_0 + k \tilde{f}(\xi) = \tilde{g}(\xi).
	\]
 Rearranging gives
	\[
	\tilde{f}(\xi)[\xi^{n}(1 - r \xi^{-m})^{l} + k]
	= \xi^{n-1}(1 - r \xi^{-m})^{l}\tau_0 + \tilde{g}(\xi).
	\]
%%--------------------------------------------	
    \[
H(\xi) := \frac{1}{\xi^{n}\big(1 - r\xi^{-m}\big)^{l} + k}.
\]

Then we have
\[
\widetilde{f}(\xi)
= H(\xi)\,\widetilde{g}(\xi)
+ \tau_{0}\,\xi^{n-1}\big(1 - r\xi^{-m}\big)^{l}H(\xi).
\]

Taking the inverse Laplace transform termwise gives
\[
f(t)= (h *g)(t)+\tau_{0}\mathcal{L}^{-1}\Big\{\frac{\xi^{n-1}(1 - r \xi^{-m})^{l}}{\xi^{n}(1 - r \xi^{-m})^{l} + k}\Big\}(t),
\]
where, $h(t) = \mathcal{L}^{-1}\{H(\xi)\}(t),
~
g(t) = \mathcal{L}^{-1}\{\widetilde{g}(\xi)\}(t),~
(h * g)(t) = \int_{0}^{t} h(t - s)\,g(s)\,ds.
$

    
    
    
    
    
    
    
    
    
    
    
    
    
    
    
    
    %%---------------------------------------
 %    Hence
	% \[
	% \tilde{f}(\xi)
	% = \frac{\xi^{n-1}(1 - r \xi^{-m})^{l}\tau_0}{\xi^{n}(1 - r \xi^{-m})^{l} + k}
	% + \frac{\tilde{g}(\xi)}{\xi^{n}(1 - r \xi^{-m})^{l} + k}.
	% \]
	% \begin{equation}
	%     \tilde{f}(\xi)
	% = \dfrac{\tau_0}{\dfrac{\xi^{n}(1 - r \xi^{-m})^{l} + k}{\xi^{n-1}(1 - r \xi^{-m})^{l}}}
	% +\frac{\tilde{g}(\xi)}{\xi^{n-ml}(\xi^m  - r )^{l} + k}. 
	% \end{equation}
 %    \begin{equation}
	%     \tilde{f}(\xi)
	% = \dfrac{\tau_0}{\xi +\dfrac{ k}{\xi^{n-1}(1 - r \xi^{-m})^{l}}}
	% + \frac{\tilde{g}(\xi)}{\dfrac{(\xi^m  - r )^{l}}{\xi^{ml-n}} + k}.
	% \end{equation}
 %    \begin{equation}
	%     \tilde{f}(\xi)
	% = \dfrac{\tau_0}{\xi +\dfrac{ k}{\xi^{n-1}\xi^{-ml} (\xi^m  - r)^{l}}}
	% + \dfrac{\tilde{g}(\xi)}{\dfrac{(\xi^m  - r )^{l}}{\xi^{ml-n}}}\left[1+\dfrac{k}{\dfrac{(\xi^m  - r )^{l}}{\xi^{ml-n}}}\right]^{-1}
	% \end{equation}
 %     \begin{equation}
	%     \tilde{f}(\xi)
	% = \dfrac{\tau_0}{\xi +\dfrac{ k\xi^{ml-n+1} }{ (\xi^m  - r)^{l}}}
	% +\tilde{g}(\xi)\frac{\xi^{ml-n}}{(\xi^m  - r )^{l}}\sum_{u=0}^\infty \left[\dfrac{-k}{\dfrac{(\xi^m  - r )^{l}}{\xi^{ml-n}}}\right]^u
	% \end{equation}
 %      \begin{equation}
	%     \tilde{f}(\xi)
	% = \dfrac{\tau_0}{\xi\left[1 +\dfrac{ k\xi^{ml-n} }{ (\xi^m  - r)^{l}}\right]}
	% + \tilde{g}(\xi)\frac{\xi^{ml-n}}{(\xi^m  - r )^{l}}\sum_{u=0}^\infty (-k)^u\left[\dfrac{\xi^{ml-n}}{(\xi^m  - r )^{l}}\right]^u
	% \end{equation}
 %    \begin{equation}
	%     \tilde{f}(\xi)
	% = \dfrac{\tau_0}{\xi}\left[1 +\dfrac{ k\xi^{ml-n} }{ (\xi^m  - r)^{l}}\right]^{-1}
	% \tilde{g}(\xi)\frac{\xi^{ml-n}}{(\xi^m  - r )^{l}}\sum_{u=0}^\infty (-k)^u\left[\dfrac{\xi^{ml-n}}{(\xi^m  - r )^{l}}\right]^u
	% \end{equation}
 %      \begin{equation}
	%     \tilde{f}(\xi)
	% = \dfrac{\tau_0}{\xi}\sum_{u=0}^{\infty}(-1)^u 
 %    \left[\dfrac{ k\xi^{ml-n} }{ (\xi^m  - r)^{l}}\right]^{u}
	% + \tilde{g}(\xi)\frac{\xi^{ml-n}}{(\xi^m  - r )^{l}}\sum_{u=0}^\infty (-k)^u\left[\dfrac{\xi^{ml-n}}{(\xi^m  - r )^{l}}\right]^u
	% \end{equation}
 %       \begin{equation}
	%     \tilde{f}(\xi)
	% = {\tau_0}\sum_{u=0}^{\infty}(-k)^u 
 %    \left[\dfrac{ \xi^{ml-n-1} }{ (\xi^m  - r)^{l}}\right]^{u}
	% + \tilde{g}(\xi)\frac{\xi^{ml-n}}{(\xi^m  - r )^{l}}\sum_{u=0}^\infty (-k)^u\left[\dfrac{\xi^{ml-n}}{(\xi^m  - r )^{l}}\right]^u
	% \end{equation}
    
 % %   first term is the Laplace transform of $t^{n}\mathbb{E}^l_{m,n+1}(rt^n)$
 % By taking the inverse Laplace transform we obtain the required result
    \end{proof}

\end{document}
